A proactive streaming assistant must decide, at every moment, whether to speak,
and that decision is almost always a scalar confidence compared against a
per-task threshold. We fit such a threshold table for Qwen2.5-Omni-7B on OmniPro Online and
report that it cannot be identified from the available data --- together with the
measurements that establish why, and the one change to the same gate that does
work. Auditing first, we find the shipped gate ignored its own threshold table
(firing on a fixed global constant) and treated the first tick of every video as a
rising edge, consuming an entire video's emission budget before any content
arrived; both are repaired before measurement. We then test each fitted threshold
by a paired bootstrap over videos that replays the actual selection rule. No task's
selection beats its best rival outside noise, selections recur in only 25--61\,\% of
resamples, and one recurs at 6\,\% against an 8\,\% chance rate. The cause is that
the gating score does not rank event-adjacent ticks above quiet ones: per-task AUC
is 0.431--0.551 with every interval covering chance, and this survives sweeping the
tick clock over $\pm 10$\,s, widening the matching tolerance, and confirming the
score is well spread over 232--284 distinct values. A threshold on such a score buys
emission volume only --- and volume has a single useful setting, which pooled over
all tasks is recovered stably in 94\,\% of resamples while nine per-task values are
not recovered at all. We argue that a per-task threshold table is a set of fitted
parameters and should be held to the corresponding standard: a resampling check
that the values recur, and a comparison against one global threshold. Both re-use
existing runs at no additional compute cost, and here each independently reverses
the conclusion the point estimates would have supported.
